\begin{table}[htbp]
\centering
\footnotesize
\caption{Problema G, Passos 2 e 4 (google/gemini-2.5-flash-lite): zero-shot e métrica em R$^2$, R$^3$ e R$^4$ por subproblema, nomes de classe e de atributo.}
\label{tab:bloco3_G_passos_2_e_4_R2_R3_R4_flashlite}
\begin{tabular}{llllrrrrrr}
\toprule
Subproblema & Classes & Atributos & Espaço & Acur.\ LLM vs real (n erros) & Acur.\ vs real (inv.\ perm.) & Fid.\ Perc. & Fid.\ NNLS & Acur.\ Perc.\ vs real (n erros) & Acur.\ NNLS vs real \\
\midrule
G1 (normalizado) & Mulher / Homem & peso/altura & R$^2$ & 35,2 (45,3) & — & 71,4 (70,0–72,9) & 71,4 & 10,5 (62,7) & 10,5 \\
G1 (normalizado) & Mulher / Homem & peso/altura & R$^3$ & 35,2 (45,3) & — & 71,4 (70,0–72,9) & 74,8 & 10,5 (62,7) & 13,8 \\
G1 (normalizado) & Mulher / Homem & peso/altura & R$^4$ & 35,2 (45,3) & — & 71,9 (71,4–72,9) & 77,1 & 11,0 (62,3) & 21,0 \\
\midrule
G1 (normalizado) & A / B & peso/altura & R$^2$ & 18,6 (57,0) & 81,4 & 93,8 (91,4–95,7) & 94,8 & 16,2 (58,7) & 18,1 \\
G1 (normalizado) & A / B & peso/altura & R$^3$ & 18,6 (57,0) & 81,4 & 93,8 (91,4–95,7) & 93,8 & 16,2 (58,7) & 20,0 \\
G1 (normalizado) & A / B & peso/altura & R$^4$ & 18,6 (57,0) & 81,4 & 91,9 (91,4–92,9) & 94,3 & 19,0 (56,7) & 18,6 \\
\midrule
G1 (normalizado) & A / B & x1/x2 & R$^2$ & 16,2 (58,7) & 83,8 & 91,0 (90,0–91,4) & 93,8 & 10,0 (63,0) & 13,8 \\
G1 (normalizado) & A / B & x1/x2 & R$^3$ & 16,2 (58,7) & 83,8 & 91,0 (90,0–91,4) & 93,3 & 10,0 (63,0) & 12,4 \\
G1 (normalizado) & A / B & x1/x2 & R$^4$ & 16,2 (58,7) & 83,8 & 91,0 (90,0–91,4) & 92,9 & 10,0 (63,0) & 12,9 \\
\midrule
\textbf{G2 (kg e m)} & \textbf{Mulher / Homem} & \textbf{peso\_kg/altura\_m} & \textbf{R$^2$} & \textbf{84,3 (11,0)} & \textbf{—} & \textbf{92,9 (91,4–94,3)} & \textbf{93,8} & \textbf{91,4 (6,0)} & \textbf{88,6} \\
\textbf{G2 (kg e m)} & \textbf{Mulher / Homem} & \textbf{peso\_kg/altura\_m} & \textbf{R$^3$} & \textbf{84,3 (11,0)} & \textbf{—} & \textbf{92,9 (91,4–94,3)} & \textbf{96,7} & \textbf{91,4 (6,0)} & \textbf{87,6} \\
\textbf{G2 (kg e m)} & \textbf{Mulher / Homem} & \textbf{peso\_kg/altura\_m} & \textbf{R$^4$} & \textbf{84,3 (11,0)} & \textbf{—} & \textbf{92,9 (91,4–94,3)} & \textbf{97,6} & \textbf{91,4 (6,0)} & \textbf{86,7} \\
\midrule
G2 (kg e m) & A / B & peso\_kg/altura\_m & R$^2$ & 46,7 (37,3) & 54,3 & 71,4 (65,7–77,1) & 72,4 & 23,8 (53,3) & 31,4 \\
G2 (kg e m) & A / B & peso\_kg/altura\_m & R$^3$ & 46,7 (37,3) & 54,3 & 73,3 (67,1–80,0) & 77,1 & 24,8 (52,7) & 36,2 \\
G2 (kg e m) & A / B & peso\_kg/altura\_m & R$^4$ & 46,7 (37,3) & 54,3 & 72,4 (64,3–81,4) & 75,7 & 21,9 (54,7) & 31,0 \\
\bottomrule
\multicolumn{10}{p{0.95\linewidth}}{\footnotesize Problema G, Passo 2 (zero-shot nas 70 amostras de treino de cada semente) e métrica estimada sobre esses rótulos (Passo 4); médias das sementes (fidelidade com mín–máx). Fid.\ = fidelidade (vs LLM) de cada estimador; Acur.\ vs real = acurácia (vs rótulo real) do LLM e da métrica, contexto, com o número médio de erros entre parênteses (contagem auxiliar). Acur.\ vs real (inv.\ perm.) = max($a$, 1 $-$ $a$) nas linhas com classes neutras A/B: em zero-shot e com nomes neutros a correspondência entre nome e classe é arbitrária e a acurácia literal pode sair invertida (a métrica, estimada dos rótulos do LLM, herda a mesma permutação). Atributos extras (x1·x2; x1², x2²) só na estimação. Combinações ausentes = o LLM colapsou no zero-shot (classe única ou minoritária com menos de 5 pontos) nas três sementes. Em negrito: G2 com nomes semânticos.} \\
\end{tabular}
\end{table}
